% Options for packages loaded elsewhere
\PassOptionsToPackage{unicode}{hyperref}
\PassOptionsToPackage{hyphens}{url}
\PassOptionsToPackage{dvipsnames,svgnames,x11names}{xcolor}
\documentclass[
  10pt,
  table]{article}
\usepackage{xcolor}
\usepackage[letterpaper,left=0.7in,right=0.7in,top=0.7in,bottom=0.75in]{geometry}
\usepackage{amsmath,amssymb}
\setcounter{secnumdepth}{-\maxdimen} % remove section numbering
\usepackage{iftex}
\ifPDFTeX
  \usepackage[T1]{fontenc}
  \usepackage[utf8]{inputenc}
  \usepackage{textcomp} % provide euro and other symbols
\else % if luatex or xetex
  \usepackage{unicode-math} % this also loads fontspec
  \defaultfontfeatures{Scale=MatchLowercase}
  \defaultfontfeatures[\rmfamily]{Ligatures=TeX,Scale=1}
\fi
\usepackage{lmodern}
\ifPDFTeX\else
  % xetex/luatex font selection
\fi
% Use upquote if available, for straight quotes in verbatim environments
\IfFileExists{upquote.sty}{\usepackage{upquote}}{}
\IfFileExists{microtype.sty}{% use microtype if available
  \usepackage[]{microtype}
  \UseMicrotypeSet[protrusion]{basicmath} % disable protrusion for tt fonts
}{}
\usepackage{setspace}
\makeatletter
\@ifundefined{KOMAClassName}{% if non-KOMA class
  \IfFileExists{parskip.sty}{%
    \usepackage{parskip}
  }{% else
    \setlength{\parindent}{0pt}
    \setlength{\parskip}{6pt plus 2pt minus 1pt}}
}{% if KOMA class
  \KOMAoptions{parskip=half}}
\makeatother
\usepackage{longtable,booktabs,array}
\usepackage{caption}
\captionsetup[table]{skip=6pt}
\newcounter{none} % for unnumbered tables
\usepackage{calc} % for calculating minipage widths
% Correct order of tables after \paragraph or \subparagraph
\usepackage{etoolbox}
\makeatletter
\patchcmd\longtable{\par}{\if@noskipsec\mbox{}\fi\par}{}{}
\makeatother
% Allow footnotes in longtable head/foot
\IfFileExists{footnotehyper.sty}{\usepackage{footnotehyper}}{\usepackage{footnote}}
\makesavenoteenv{longtable}
\usepackage{graphicx}
\makeatletter
\newsavebox\pandoc@box
\newcommand*\pandocbounded[1]{% scales image to fit in text height/width
  \sbox\pandoc@box{#1}%
  \Gscale@div\@tempa{\textheight}{\dimexpr\ht\pandoc@box+\dp\pandoc@box\relax}%
  \Gscale@div\@tempb{\linewidth}{\wd\pandoc@box}%
  \ifdim\@tempb\p@<\@tempa\p@\let\@tempa\@tempb\fi% select the smaller of both
  \ifdim\@tempa\p@<\p@\scalebox{\@tempa}{\usebox\pandoc@box}%
  \else\usebox{\pandoc@box}%
  \fi%
}
% Set default figure placement to htbp
\def\fps@figure{htbp}
\makeatother
\setlength{\emergencystretch}{3em} % prevent overfull lines
\providecommand{\tightlist}{%
  \setlength{\itemsep}{0pt}\setlength{\parskip}{0pt}}
% Phase 2 look, rendered through LaTeX (XeTeX via tectonic)
\usepackage{fontspec}
\setmainfont{Helvetica Neue}[Ligatures=TeX,BoldFont={Helvetica Neue Bold},ItalicFont={Helvetica Neue Italic},BoldItalicFont={Helvetica Neue Bold Italic}]
\setsansfont{Helvetica Neue}[Ligatures=TeX,BoldFont={Helvetica Neue Bold},ItalicFont={Helvetica Neue Italic},BoldItalicFont={Helvetica Neue Bold Italic}]
\newfontfamily\fallbackfont{Arial Unicode MS}
\usepackage{newunicodechar}
\newunicodechar{₹}{Rs.\,}
\newunicodechar{→}{{\fallbackfont →}}
\newunicodechar{≈}{{\fallbackfont ≈}}
\newunicodechar{≥}{{\fallbackfont ≥}}
\newunicodechar{≤}{{\fallbackfont ≤}}
\newunicodechar{×}{{\fallbackfont ×}}
\newunicodechar{✓}{{\fallbackfont ✓}}

\definecolor{pblue}{HTML}{0D47A1}
\definecolor{pblue2}{HTML}{1565C0}
\definecolor{pslate}{HTML}{37474F}
\definecolor{prule}{HTML}{C8D8F0}
\definecolor{pzebra}{HTML}{F4F7FB}
\definecolor{pline}{HTML}{DDE3EA}
\definecolor{ptext}{HTML}{232323}

\usepackage{etoolbox}
\usepackage{titlesec}
\usepackage{titletoc}
\usepackage{enumitem}
\usepackage{caption}
\usepackage{float}
\usepackage{adjustbox}
\usepackage{booktabs}
\usepackage{tabularx}
\usepackage{array}
\usepackage{fancyhdr}

\AtBeginDocument{\color{ptext}}
\setlength{\parindent}{0pt}
\setlength{\parskip}{5.5pt plus 1pt}
\setlist{itemsep=1.5pt,topsep=2pt,parsep=0pt,leftmargin=*}
\renewcommand{\tightlist}{\setlength{\itemsep}{1.5pt}\setlength{\parskip}{0pt}}

\titleformat{\section}{\color{pblue}\Large\bfseries}{}{0pt}{}[{\vspace{-1pt}\color{prule}\titlerule[1.3pt]}]
\titlespacing*{\section}{0pt}{16pt}{7pt}
\titleformat{\subsection}{\color{pblue2}\large\bfseries}{}{0pt}{}
\titlespacing*{\subsection}{0pt}{11pt}{4pt}
\titleformat{\subsubsection}{\color{pslate}\normalsize\bfseries}{}{0pt}{}
\titlespacing*{\subsubsection}{0pt}{8pt}{3pt}
\titleformat{\paragraph}[runin]{\color{pslate}\normalsize\bfseries}{}{0pt}{}
\titlespacing*{\paragraph}{0pt}{6pt}{6pt}

\titlecontents{section}[0pt]{\addvspace{5pt}\bfseries\color{pblue}}{}{}{\hfill\contentspage}
\titlecontents{subsection}[1.2em]{\small}{}{}{\titlerule*[3pt]{.}\contentspage}

\captionsetup{font={small,it},labelfont={bf,color=pblue2},labelsep=period,skip=4pt,justification=raggedright,singlelinecheck=false}
\floatplacement{figure}{H}
\setcounter{secnumdepth}{0}

% tables: header row blue, zebra body, small type
\AtBeginEnvironment{longtable}{\footnotesize\renewcommand{\arraystretch}{1.18}\rowcolors{2}{white}{pzebra}}
\setlength{\LTpre}{6pt}\setlength{\LTpost}{8pt}
\setlength{\tabcolsep}{4.2pt}
\providecommand{\tblhead}[1]{\cellcolor{pblue}\textcolor{white}{\textbf{#1}}}

\pagestyle{fancy}
\fancyhf{}
\renewcommand{\headrulewidth}{0pt}
\fancyfoot[C]{\footnotesize\color{pslate}\thepage}
\fancypagestyle{plain}{\fancyhf{}\renewcommand{\headrulewidth}{0pt}\fancyfoot[C]{\footnotesize\color{pslate}\thepage}}

% title block, as in the Phase 2 documents
\gdef\psubtitle{}
\providecommand{\subtitle}[1]{}
\renewcommand{\subtitle}[1]{\gdef\psubtitle{#1}}
\makeatletter
\renewcommand{\maketitle}{%
  \begingroup\raggedright
  {\color{pblue}\fontsize{24}{27}\selectfont\bfseries \@title\par}
  \vspace{3pt}{\color{pblue}\rule{\linewidth}{2.4pt}}\par\vspace{5pt}
  {\color{pslate}\large \psubtitle\par}
  \vspace{4pt}
  \endgroup}
\makeatother
\usepackage{bookmark}
\IfFileExists{xurl.sty}{\usepackage{xurl}}{} % add URL line breaks if available
\urlstyle{same}
% fallback for those not using the hyperref driver hyperxmp:
\makeatletter
\@ifundefined{xmpquote}{\newcommand{\xmpquote}[1]{#1}}{}
\makeatother
\hypersetup{
  pdftitle={System Context Brief},
  colorlinks=true,
  linkcolor={pblue2},
  filecolor={Maroon},
  citecolor={Blue},
  urlcolor={pblue2},
  pdfcreator={LaTeX via pandoc}}

\title{System Context Brief}
\usepackage{etoolbox}
\makeatletter
\providecommand{\subtitle}[1]{% add subtitle to \maketitle
  \apptocmd{\@title}{\par {\large #1 \par}}{}{}
}
\makeatother
\subtitle{IIIT Hyderabad Breakfast Mess System}
\author{}
\date{}

\begin{document}
\maketitle

\renewcommand*\contentsname{Contents}
{
\hypersetup{linkcolor=pslate}
\setcounter{tocdepth}{2}
\tableofcontents
}
\setstretch{1.08}
\section{System Boundary and Context}\label{system-boundary-and-context}

\subsection{System Boundary}\label{system-boundary}

The system under study is the breakfast meal as it is registered, billed, cooked and served in the four residential dining halls at IIIT Hyderabad: Kadamba, Palash, Bakul and Yuktāhār. Three layers sit inside the boundary.

The registration and billing layer is the dining portal (dining.iiit.ac.in), where students book, cancel and pay for meals, together with the QR scan at the counter that records who ate. The kitchen layer is the outside caterer that runs each hall: Prism at Kadamba, ABC Hospitality Services at Yuktāhār, and Vijayalakshmi Caterers at Bakul, which cooks at its Hafizpet kitchen (single-sourced). It covers quantity planning, cooking, food-safety checks before service, the counter, top-up during service, and what happens to leftovers. Food for Bakul and Palash comes from one approved off-site kitchen (confirmed), so for those two halls the layer also includes the transport link and the hot holding of food at the hall. Whether Vijayalakshmi Caterers also supplies Palash is a candidate. The governance layer is the set of bodies that fix the menu, the registration and billing rules and the operator contracts.

Several parties sit on the boundary, exchanging food, money or information with the system without being run by it: the operators\textquotesingle{} suppliers, the garbage contractor, the institution\textquotesingle s site team (which relays complaints) and facility team (weekly quality inspection), and the tender through which each operator is hired. People outside the registration also eat from the same pots. At Kadamba these are 30 to 50 outsiders such as night-shift security, plus about half of the operator\textquotesingle s own staff (single-sourced). The academic class schedule sits outside the boundary but shapes it, because the two run on overlapping, independently set timings. On-campus alternatives (Vindhya Canteen, the juice canteens, delivery apps) sit at the outer edge: a student who skips the mess eats from one of them.

Lunch and dinner at the same halls fall outside the boundary and are the comparison that justifies it. The same portal, rules and operator produce very different turnout at the two meals. At Yuktāhār in September 2026, lunch attendance ran from 56 to 89 percent of registrations (15 of 20 days at 70 percent or above) while breakfast on the same days ran from 18 to 40 percent (confirmed, from the operator\textquotesingle s attendance register). The problem is specific to breakfast.

\begin{figure}
\centering
\includegraphics[width=1\linewidth,height=\textheight,keepaspectratio,alt={System boundary: three layers inside, boundary actors on the ring, class schedule and alternatives outside.}]{Assets/diagram1-system-boundary.png}
\caption{System boundary: three layers inside, boundary actors on the ring, class schedule and alternatives outside.}
\end{figure}

\subsection{Context}\label{context}

The published breakfast window is 7:30 to 9:30 AM. At Kadamba in April 2026 the first scans came at about 7:28, service closed between 9:35 and 9:45, and the latest scan was at 10:07 (confirmed, scan records). Kitchen staff start between 5:00 and 5:30 AM. Batch items are ready by about 7:30, and at Kadamba live items such as dosa are made at the counter from 7:00 (single-sourced). Bakul\textquotesingle s food is cooked off site, so labour-heavy items such as aloo paratha may be started around 3 AM to reach the hall by about 7 (single-sourced).

On one day in September the platform showed breakfast capacity of 2,755 registrations: Kadamba 700 vegetarian and 600 non-vegetarian, Bakul 350 and 350, Palash 400, Yuktāhār 340 regular and 15 Jain. Registered prices are ₹48 (vegetarian) and ₹66 (non-vegetarian) at Kadamba, Bakul and Palash, and ₹53 at Yuktāhār. Without prior registration the same breakfast costs ₹81, ₹111 and ₹92, about 70 percent more (confirmed, posted rates).

Registration is high and flat; turnout is low and variable. Kadamba took 32,162 breakfast registrations in April 2026, about 1,070 a day on every day of the week including Sunday, and scanned 11,941 breakfasts: 37.1 percent turnout, 62.9 percent no-shows (confirmed, scan records). Yuktāhār\textquotesingle s breakfast registrations ran from 228 to 404 a day in September, with turnout averaging about 28 percent (confirmed, operator\textquotesingle s register). These are the only two halls with turnout records. Bakul\textquotesingle s manager estimates that at most about 60 percent of registered students come (single-sourced; nothing written behind it).

Registration is automatic. Every student is registered for every meal by default, on campus or not, holidays included. The only exit is a cancellation, allowed five times per meal type per month and at least four days ahead; beyond that a meal cannot be removed from the bill (confirmed). Billing is monthly and follows registration, so a student pays for every meal they are registered for, eaten or not. Operators describe counts reaching them in weekly blocks, while students describe registering monthly. How the default is batched on the portal is unconfirmed (single-sourced).

The platform is shared but the kitchens are not. Kadamba serves South Indian food, vegetarian and non-vegetarian planned separately, from an on-site kitchen run by Prism. Bakul serves North Indian vegetarian and non-vegetarian food and Palash North Indian vegetarian food, both from the off-site kitchen until the Felicity Kitchen and Dining Facility opens, expected December 2026 (confirmed). Yuktāhār serves Satvik vegetarian food with a Jain variant, from an on-site kitchen run by ABC Hospitality Services (confirmed). Bakul\textquotesingle s food travels about 8 km from Hafizpet, roughly 20 minutes early in the morning and 40 in normal traffic, in insulated boxes holding about 60 to 72 °C, and stands in a hot-water setup at the hall (single-sourced). Each operator buys its own ingredients on its own cadence. The four halls are unified only at the registration, billing and access layer. (Revised from the first submission, which had Bakul\textquotesingle s vegetarian food cooked at Palash and carried over.)

Five adjacent systems act on the mess without being controlled by it.

\begin{itemize}
\tightlist
\item
  \textbf{Academic administration.} Classes run 8:30 AM to 6:40 PM with a 1:00 to 2:00 PM break, and most students have few classes on a given day. No channel links the mess to the academic calendar. Exams, events and holidays reach a kitchen only if someone raises them in its planning meeting: "we don\textquotesingle t come to know until we sit in the meeting," as the Yuktāhār site head put it (confirmed).
\item
  \textbf{Institutional governance.} A Mess Committee sets the menu and mess policy with student input about a month ahead; a Mess Office runs the portal, from which operators read their counts; the Warden holds vendor management and structural change. The Mess Committee is reported to be faculty-chaired, and a policy change tightening cancellation rules is reported to have been made to cut staff workload, both from a press summary we could not read in full (candidate). Operators speak of a "CDS team." The Kadamba operator describes its line of authority as CFS, then CDS, then Prism, then an operations manager, a manager, a supervisor and the service staff (single-sourced). Who heads CDS is open: the operator names one person and an earlier institutional interview names another. Whether CDS is the Mess Office, and how CFS relates to the Committee and the Warden, is unconfirmed.
\item
  \textbf{The operators and their contracts.} Each firm is chosen by tender and paid monthly as a lump sum on plates, with deductions (single-sourced). Whether "plates" means registered or served decides who bears the cost of a no-show, and the evidence disagrees. Kadamba\textquotesingle s operator talks of being paid for what is registered, and Bakul\textquotesingle s says its payment is on plates served (both single-sourced). The contracts may differ, or one remark may describe a duty rather than the payment formula. The basis is unconfirmed for every operator and Bakul\textquotesingle s answer is not applied to the others. The institution sets the menu and the rules; the operator decides how much to cook.
\item
  \textbf{Mess Cell.} A student WhatsApp group where people who hold a breakfast registration they will not use sell it at a negotiated price, transferring the personal QR credential. It runs outside the portal. Kadamba\textquotesingle s manager describes students registering "even if you are in the village" while a friend eats in their place, which corroborates informal transfer (single-sourced).
\item
  \textbf{The wider campus food ecosystem.} Vindhya Canteen (Monday to Saturday from 10:00 AM), the juice canteens (mornings) and delivery apps (orders until 11:00 PM, riders kept outside the gate) all compete for the same students.
\end{itemize}

\begin{figure}
\centering
\includegraphics[width=1\linewidth,height=\textheight,keepaspectratio,alt={Context map: who administers, cooks, eats and absorbs surplus around the registration count.}]{Assets/diagram2-context-map.png}
\caption{Context map: who administers, cooks, eats and absorbs surplus around the registration count.}
\end{figure}

\section{Purpose and Intended Outcomes}\label{purpose-and-intended-outcomes}

\subsection{Purpose}\label{purpose}

The system exists to give resident students reliable, adequate, good-quality food each morning at a predictable, affordable cost, covering vegetarian, non-vegetarian and Jain requirements, in support of health, energy and academic performance. The institution states its objective as to "provide nutritious, hygienic, and affordable meals while ensuring efficient operations and minimizing food wastage through advance meal planning." Advance registration is the mechanism it relies on to match food to demand.

\subsection{Actual Mechanism}\label{actual-mechanism}

Registration is a default, not a choice. The only way out is one of five cancellations per meal type per month, made at least four days ahead (confirmed). The booking therefore records the absence of an opt-out. It stays flat across the week: Kadamba\textquotesingle s roughly 1,070 daily registrations hold from Monday to Sunday while turnout runs from 42.1 percent on Tuesdays to 27.3 percent on Sundays (confirmed, scan records). Billing follows the booking, not the meal.

The kitchen does not cook to the booking. Each operator reads the closed count ahead of service, the day before, and applies its own discount.

Yuktāhār plans in daily meetings of all staff, one for the next day\textquotesingle s breakfast and lunch and one later for snacks and dinner. The inputs are portal registrations, the Wastage Book entry for the matching day of the menu cycle (the menu repeats fortnightly, so a week-3 Thursday is planned from the week-1 Thursday), an attendance register grouped by weekday, the month\textquotesingle s turnout pattern, and any holidays or events someone mentions. The Wastage Book records, per item, what was issued, cooked and left, the percentage wasted and the time the item ran out, and it works as a rule: an item that runs out well before close is re-cooked, one that runs out near close is not (confirmed, operator and records). During service the first half hour is used as a sample to project the rest, which drives top-up cooking (single-sourced). That sample is taken before the closing crest, in which about a quarter of Kadamba\textquotesingle s diners arrive; whether Yuktāhār has the same crest is unmeasured (candidate).

Kadamba cooks a first batch for about 70 percent of registrations, planned separately for vegetarian and non-vegetarian and portioned by a gram standard per dish (four idlis or 100 g of sambar a head, for example). The rest of the raw material stays ready, and a supervisor watching the counters tells the kitchen how much more to cook, ready in 10 to 15 minutes. No record carries one day\textquotesingle s outcome into the next day\textquotesingle s ratio, and the manager says of past attempts that "it was wasted many times" (single-sourced). Biryani lunch, outside the boundary, is cooked at 100 percent of registrations (single-sourced).

Bakul cooks 70 to 80 percent of registrations, the share depending on the item. The manager sets it from past attendance, reviewed weekly: if Mondays have run near 70 percent he cooks for about 80 percent that day. Whether the history is written down or carried in his head is not known. His rule is that there should be no shortage even if some food is wasted, so on low weekends he still keeps food for about 70 percent of registrations. During service he checks every 30 minutes and phones the Hafizpet kitchen with a rough quantity, which is cooked there and driven over on a lead time nobody has stated (single-sourced). (Revised from the first submission, which treated Bakul\textquotesingle s ratio as fixed like Kadamba\textquotesingle s.)

Against the records, these rules sit well above breakfast turnout. A 70 percent first batch at Kadamba is about twice the 37.1 percent who arrive (candidate). On the Bakul manager\textquotesingle s own figures, his first batch and his weekend floor both exceed the highest turnout he expects, about 60 percent, and the margin is larger if Bakul resembles the two recorded halls (candidate; no Bakul records exist). Yuktāhār\textquotesingle s site head believes breakfast turnout is about 50 percent, while his own register averages about 28 (confirmed). Even the operator with the fullest books misremembers its own number.

Surplus has four exits. Yuktāhār reuses leftovers: idli becomes idli upma or idli poha for lunch, batter goes back to the fridge, surplus milk is set as curd (confirmed). Kadamba\textquotesingle s surplus feeds the 30 to 50 outside diners, for whom about 30 extra portions are cooked, and about half the operator\textquotesingle s staff (single-sourced); the team also saw food cooked ahead, reheated and moved between vessels several times before it is eaten, which suggests some surplus is held into later service (candidate). The rest goes to the bin. The earlier finding that composting absorbs the surplus is withdrawn: Kadamba\textquotesingle s waste goes to a garbage contractor with no composting (single-sourced), and whether the contractor takes it to the compost facility the institution describes is unverified. The earlier figure of 20 to 30 staff eating leftovers is replaced by the outside-diner figure above.

Kadamba\textquotesingle s operator puts daily waste at 5 to 10 percent, without saying of what (single-sourced). A 70 percent first batch against 37 percent turnout could leave that little waste only if the outside diners, reheating and small top-ups absorb most of the difference, or if the figure is measured on another base. Which holds is unknown, and the earlier kilogram figures for Kadamba waste are withdrawn.

The operators read the same count with different aims. Kadamba treats it as its obligation: it cooks to registrations whether or not people want to eat, "because it is a responsibility and they are getting paid for it." Yuktāhār names best quality, less waste, lower operating cost and higher profit together as its goals, so waste is a cost to it. Bakul states a third aim, no shortage at the price of some waste (single-sourced for all three).

\subsection{Gap Between Purpose and Mechanism}\label{gap-between-purpose-and-mechanism}

Advance planning works only if advance registration expresses an intention to eat. Under automatic registration it does not, because the count records who has not opted out. For breakfast the distortion is large: registrations overstate turnout about 2.7 times at Kadamba and about 3.5 times at Yuktāhār (confirmed, records). Every operator plans against a number it knows is inflated and discounts it by its own rule. The gap is already measured, in the institution\textquotesingle s scan records and in the operators\textquotesingle{} registers, but nothing carries that measurement into the registration or billing rules, and inside the kitchens it has not produced a discount matched to breakfast turnout. At Bakul it is judged weekly by one manager and may exist only in his head.

\subsection{Intended Outcomes}\label{intended-outcomes}

\begin{itemize}
\tightlist
\item
  \textbf{Every registered student gets a breakfast matching their dietary requirement inside the window.} Under strain at the close. Kadamba reopens the door for 5 to 10 minutes for latecomers, and students arriving after 9:30 also come in by the back door and are served because the mess cannot refuse them (single-sourced). Kadamba and Yuktāhār cook more during the final surge (single-sourced and confirmed). Bakul keeps food for about 50 people at 9:20 against a late tail its manager puts at 10 to 15 arrivals in the last ten minutes (single-sourced).
\item
  \textbf{Quantity and quality stay consistent with actual demand, avoiding both shortage and waste.} Not met, and the best-measured failure. Kadamba\textquotesingle s first batch is set well above turnout, yet it runs short in the last fifteen minutes when arrivals peak (turnout and arrival peak confirmed from records; ratios and shortages single-sourced). Quality is also at risk. At Kadamba the team judged that reheated, transferred food loses quality (candidate). At Bakul food arrives hot but steam in the boxes softens items such as aloo paratha: the transport keeps temperature and can cost texture (single-sourced).
\item
  \textbf{Costs stay affordable and predictable across a semester.} Predictable, not efficient. In April 2026, 20,221 registered breakfasts at Kadamba were billed and not eaten (confirmed, scan records).
\item
  \textbf{Access does not depend on navigating informal workarounds.} Not met. Students resell registrations outside the system, and operators describe friends eating on another student\textquotesingle s registration.
\item
  \textbf{The system responds to changing student needs through its own governance.} Met inside each operator, not between operator and institution. Operators adjust quantities and item mix from what they see, but the institution\textquotesingle s exam, event and holiday calendar reaches them only informally, and they have little say over the menu.
\end{itemize}

\section{Visible Symptoms and Underlying Issues}\label{visible-symptoms-and-underlying-issues}

{\def\LTcaptype{none} % do not increment counter
\begin{longtable}[]{@{}
  >{\raggedright\arraybackslash}p{(\linewidth - 2\tabcolsep) * \real{0.4595}}
  >{\raggedright\arraybackslash}p{(\linewidth - 2\tabcolsep) * \real{0.5405}}@{}}
\toprule\noalign{}
\begin{minipage}[b]{\linewidth}\raggedright
\cellcolor{pblue}\color{white}\bfseries Visible Symptom
\end{minipage} & \begin{minipage}[b]{\linewidth}\raggedright
\cellcolor{pblue}\color{white}\bfseries Underlying Issue
\end{minipage} \\
\midrule\noalign{}
\endhead
\bottomrule\noalign{}
\endlastfoot
About two thirds of breakfast registrations go unused: 37.1 percent of 32,162 were eaten at Kadamba in April, about 28 percent at Yuktāhār in September. Yuktāhār lunch on the same days ran 56 to 89 percent (confirmed, records). & Registration is an automatic default, not a record of the morning\textquotesingle s decision to eat (confirmed). The same platform and rules work at Yuktāhār lunch, the one lunch with records, so the cause is specific to breakfast. Timing against the 8:30 class start, late sleep and a four-day-ahead exit are the candidate causes. (Revises the first submission\textquotesingle s working gap of about 30 percent, which fits lunch, not breakfast.) \\
Kadamba\textquotesingle s registrations stay at about 1,070 a day on every weekday including Sunday while turnout runs from 42.1 percent (Tuesday) to 27.3 percent (Sunday). Yuktāhār\textquotesingle s Sunday breakfasts run 18 to 24 percent, among its lowest. Bakul\textquotesingle s manager calls weekend turnout very low (single-sourced). & The count reflects a default set days earlier, with an exit that closes four days ahead, so it cannot follow day-to-day variation in who eats (confirmed). \\
A student registers, does not attend, and is still charged the full price. & Billing follows registration, so the cost is incurred when the default is left standing, whether or not the meal is eaten. \\
The portal caps cancellations at five per meal type per month, each at least four days ahead. Beyond five there is no way off the bill (confirmed). & The cap limits exits from a default, not bookings. An explicit limit on leaving suggests unused registrations were frequent enough to need one. \\
"Skip Meal" lets a student flag a missed meal, but the charge stands. & The tool exists to inform the kitchen and returns nothing to the student. Whether it changes what Kadamba cooks is unresolved: the institution\textquotesingle s account says declarations reach the kitchen, Kadamba\textquotesingle s operator says it cooks on registrations alone, and no operator names Skip Meal as a planning input (single-sourced). \\
Students resell unused registrations in a WhatsApp group, "Mess Cell," at a negotiated price. & Resale returns value and Skip Meal returns none, so for a student who knows in advance, the informal market beats the official tool. \\
Unused registrations persist although cancellation and resale both exist. & Both need advance knowledge. A decision made on the morning (oversleeping, a change of plan) has no mechanism at all. \\
Walk-in breakfast costs about 70 percent more (₹81 against ₹48). At Yuktāhār walk-ins are rare: unregistered, paying and faculty or staff diners together number 0 to 12 per breakfast (confirmed, records). & The price gap gives a standing reason to register even when attendance is uncertain, which inflates counts independent of intention. \\
Registration relative to capacity differed by cuisine on one September day: 51 to 79 percent at Kadamba\textquotesingle s lines, 11 to 36.5 percent at the North Indian lines (Bakul and Palash). & These are registrations as a share of capacity from one portal snapshot, which disagrees with Yuktāhār\textquotesingle s own register for the same day (116 registered on the portal, 276 in the register), so it was probably not a final count. Cuisine may influence where students register (candidate); it does not explain turnout. Kadamba\textquotesingle s turnout (37.1 percent) is only modestly above Yuktāhār\textquotesingle s (about 28), and within Kadamba the most-booked line, vegetarian, has the lowest turnout (32.4 percent against 48.5 for non-vegetarian). No records exist for Bakul or Palash. (Revises the first submission, which read these figures as attendance.) \\
Arrivals pile up at the end of the window. At Kadamba 26.5 percent of scans fall between 9:15 and 9:30 and 13.5 percent between 9:00 and 9:15, climbing to about nine a minute at 9:29 (confirmed, records). Yuktāhār\textquotesingle s operator also describes a closing rush. Kadamba plans for about 150 arrivals in the last ten minutes and serves students after close. Bakul\textquotesingle s manager describes an earlier peak, 8:00 to 8:30, with 10 to 15 in the last ten minutes (single-sourced). & Students who come arrive as late as the window allows; late waking is the candidate driver, and the one reason students most often give. Operators must keep food until the last minute, which drives both cooking "just in case" and shortage at the close. It is the one symptom where shortage and waste occur together. The curve differs by hall: a closing crest is recorded at Kadamba and described at Yuktāhār, but not at Bakul. (Revises the first submission\textquotesingle s claim that all three describe a closing rush.) \\
Kitchens cook well above turnout: Kadamba 70 percent of registrations against 37 measured; Bakul 70 to 80 percent, and about 70 on weekends, against its manager\textquotesingle s own ceiling of 60; Yuktāhār\textquotesingle s head believes 50 percent against a recorded 28. & Each operator applies its own factor to a count it knows is inflated: a fixed ratio with reactive top-up at Kadamba, a weekly judgement with a shortage floor at Bakul, a book system with first-half-hour projection at Yuktāhār. Only Yuktāhār feeds recorded outcomes back into planning. Bakul\textquotesingle s 60 percent is a belief; if Bakul resembles the recorded halls at 28 to 37 percent, it is an upper bound (candidate). \\
At Kadamba, food is cooked ahead, reheated and moved between vessels, and the team judged its quality to fall fast (candidate). & A first batch far above turnout is held, not discarded, so surplus returns as reheated food. Whether this lowers attendance is not evidenced; it is an open chain, not a loop. \\
Some items are wasted every time they appear (at Kadamba uttapam, set dosa, upma) while others move beyond the per-head standard (puri, bonda, vada, idli). A student complained that dinner sambar was dropped while charges rose. Bakul\textquotesingle s manager names pav bhaji as a draw, muesli and Chocos as items that go, and objects to some pairings students choose, such as Raj Kachori listed with chana where he would serve aloo curry (single-sourced; untested with Bakul residents). & Students and the committee choose the menu, and only operators see item-level waste and uptake. Nothing connects the two. Yuktāhār\textquotesingle s records already work item by item on a fortnightly cycle (confirmed) but inform quantities, not the menu. \\
No student can name an instance of feedback changing registration, billing or hours. & Three channels exist (an in-app rating, a public email address, student-press coverage) and none has a confirmed instance of changing those rules. The one reported change driven by students, the fortnightly menu rotation, is single-sourced. Change does happen inside operators: Yuktāhār began weighing waste in August of last year, made it blame-free, and staff reported waste falling by December; Kadamba trimmed ragi idli and raised plain idli from what sold (single-sourced). None of it reaches institutional policy. The clash between breakfast hours and classes has never been formally raised because no path exists for a cross-authority issue. \\
Breakfast (7:30 to 9:30) overlaps the 8:30 class start, and separate authorities set the two. A late-waking student says they choose class over breakfast because of the attendance policy. & The mess bodies control operations, the academic office controls the timetable, neither controls the other\textquotesingle s schedule, and nothing coordinates them. The conflict can persist because neither side can resolve it alone. \\
Students are registered for every meal whether or not they are on campus (confirmed). The portal also has a setting to notify a student of a randomly allocated meal. & The count records no opt-out. How random allocation works alongside the default is unconfirmed. (Revises the first submission, which treated only some registrations as system-generated.) \\
Some students skip for reasons unrelated to the mess: diet, sleep schedule, habit. & Non-attendance has more than one cause. Part of it would persist whatever changed in hours, pricing or process. \\
Students say attendance is more likely when a friend is going. & For some students breakfast is socially coordinated, so a registration alone does not predict who turns up. \\
\end{longtable}
}

\section{Key Actors, Institutions, Processes, Resources and Constraints}\label{key-actors-institutions-processes-resources-and-constraints}

\subsection{Actors}\label{actors}

\begin{itemize}
\tightlist
\item
  \textbf{Students.} Registrants and eaters; attendance and diet vary widely.
\item
  \textbf{Mess Committee.} A policy body that sets the menu and mess policy with student input; reported to be faculty-chaired (candidate).
\item
  \textbf{Mess Office.} The operational arm that runs the portal and issues operational policy. Operators say "CDS team" and describe a site team that relays complaints; whether this is the Mess Office is unconfirmed.
\item
  \textbf{CFS and CDS.} The institutional line above the vendor as the Kadamba operator describes it (single-sourced). Their relation to the Committee, Office and Warden is unconfirmed, and the identity of the CDS head is open.
\item
  \textbf{Warden.} Holds vendor management and structural change.
\item
  \textbf{Catering operators.} ABC Hospitality Services at Yuktāhār (confirmed); Prism, hired by tender, which runs all Kadamba operations (single-sourced); Vijayalakshmi Caterers at Bakul, from Hafizpet (single-sourced). Whether it also serves Palash is a candidate. (Revises two first-submission readings: Prism as the name of Kadamba\textquotesingle s floor team, and Vijayalakshmi\textquotesingle s hall as unnamed.)
\item
  \textbf{Bakul staff.} A manager, a month in the post but with about 20 years of North Indian food experience, who sets the first batch, checks the counter every 30 minutes and phones Hafizpet for more (single-sourced).
\item
  \textbf{Kadamba staff.} A chain under Prism of operations manager, manager, supervisor and service staff, about ten chefs under a main vegetarian chef; the supervisor watches the counters and calls for more cooking (single-sourced).
\item
  \textbf{Yuktāhār staff.} Sixteen people, no hierarchy, rotating across every job, in tasting, cutting and cleaning teams who cross-check one another\textquotesingle s records. The mess in-charge keeps the accounts and the Wastage Book, and a colleague cross-verifies the grocery register. Two project managers, with some site staff sharing the work, run records and hygiene. There is no leave policy: absence is raised in the daily meeting and colleagues agree to cover (single-sourced). (Revises an earlier five-person project team.)
\item
  \textbf{Non-student diners at Kadamba.} 30 to 50 outsiders such as night-shift security, and about half the operator\textquotesingle s staff (single-sourced). (Revises an earlier 20 to 30 staff.)
\item
  \textbf{Suppliers and the waste contractor.} Grocery, vegetable, milk and paneer suppliers for each operator, and the contractor that removes Kadamba\textquotesingle s food waste.
\item
  \textbf{Quality control.} A weekly QC inspection; at Kadamba a facility team tastes each meal with two Prism staff before service (single-sourced; the tasting was observed). Faculty visitors are also described as inspecting operator stores (single-sourced).
\item
  \textbf{Academic administration.} Sets the timetable; no authority over the mess.
\item
  \textbf{Mess Cell sellers.} A subset of registered students.
\item
  \textbf{Guests and walk-ins.} Parents, relatives, faculty, staff and other campus affiliates paying the spot rate.
\end{itemize}

\subsection{Institutions}\label{institutions}

\begin{itemize}
\tightlist
\item
  Three governance actors hold different parts of mess governance: the Mess Committee (policy), the Mess Office (operations) and the Warden (vendors and structural change). The Kadamba operator names a different line, CFS then CDS (single-sourced). How the two descriptions map onto each other is unconfirmed.
\item
  The academic administration and timetable authority, independent of all of them.
\item
  The tender and monthly per-plate payment through which each operator is contracted (single-sourced; the plate basis is unknown).
\item
  The operators\textquotesingle{} record systems. Yuktāhār keeps a set of paper books, all photographed (confirmed): a Wastage Book (per dish: issued, cooked, left, percentage, run-out time, re-cook quantity, registered and eaten counts, student plate waste), a weekday attendance register, a wastage board kept in view, a grocery issue book headed with the day\textquotesingle s registrations, a store and ingredient book, a stock register, a vegetable book and indent with spoilage remarks, a milk, fruit, gas, electricity and water register, and fridge temperature and cleaning logs. The issue book must reconcile with what was cooked, and a mismatch is raised in that day\textquotesingle s meeting (single-sourced); on the day checked, the issue book and stock register agree (confirmed). The site head describes cross-checking as "not to find fault," and weighing was made blame-free on purpose (single-sourced). Together these hold the system\textquotesingle s most detailed demand data.
\item
  The food-safety regime: an external "Eat Right Campus" five-star certification and laboratory testing of suspect food. At Kadamba every meal passes a fixed sequence: cooking temperature checked, a sample kept for 72 hours so any illness can be traced, a wrapped display plate set at the entrance, serving temperature taken per item, then a tasting by the facility team with two Prism staff and corrections. QC inspects weekly, and water TDS is read daily with a failure triggering a full check (single-sourced; display plate and tasting confirmed by observation).
\end{itemize}

\begin{figure}
\centering
\includegraphics[width=1\linewidth,height=\textheight,keepaspectratio,alt={What the dining halls post: the main overview poster, the student handbook, the food-card guide and the food-safety certificate.}]{Assets/plate-posters.jpg}
\caption{What the dining halls post: the main overview poster, the student handbook, the food-card guide and the food-safety certificate.}
\end{figure}

\subsection{Processes}\label{processes}

\begin{itemize}
\tightlist
\item
  \textbf{Registration.} Students are registered automatically for their hall and diet category, on campus or not. The student\textquotesingle s action is to opt out (confirmed).
\item
  \textbf{Billing.} Monthly, calculated from registrations, independent of attendance.
\item
  \textbf{Cancellation.} At least four days ahead, five per meal type per month, removing the charge. No other route removes a meal from the bill (confirmed).
\item
  \textbf{Skip Meal.} A flag to the kitchen that does not affect billing; its effect on Kadamba\textquotesingle s cooking is unresolved (see symptoms).
\item
  \textbf{Menu setting.} The Mess Committee sets the menu with student input; operators describe a student committee choosing from options they offer. Bakul\textquotesingle s manager says students choose combinations from a balanced-diet, protein and preference view (he mentions a target of about 30 to 40 g of protein per meal), while he would pair dishes by traditional and regional pattern and by what people eat; his account of the students\textquotesingle{} reasoning is untested (single-sourced). The menu runs on a two-week cycle, weeks 1 and 3 repeating and weeks 2 and 4 repeating (confirmed, Yuktāhār records).
\item
  \textbf{Procurement.} Each operator buys on its own cadence. Yuktāhār buys groceries fortnightly, vegetables two or three times a week, milk daily (confirmed). Kadamba buys dry goods a week ahead, some vegetables two or three days ahead, and milk, fruit, meat, eggs and most vegetables the day before (single-sourced). Yuktāhār buys from many local shops and godowns, and its store cleans and sorts raw material before issue against an expected impurity norm per item, above which the supplier is challenged or changed (cleaning confirmed, norm single-sourced); the operator visits mills to check how inputs are made (single-sourced).
\item
  \textbf{Quantity planning.} The operator reads the closed count the day before and applies its own conversion: Yuktāhār through the meeting and books, producing a written per-meal issue list (confirmed for books and list); Kadamba a fixed 70 percent with gram standards; Bakul 70 to 80 percent by item revised weekly (single-sourced). (Revises the first submission, which had the Mess Office placing food orders four days ahead. The four-day lead applies to registration; the operator sets quantities the day before.)
\item
  \textbf{Cooking and top-up.} The first batch is cooked from 5:30 (from about 3 AM for labour-heavy items bound for Bakul) and at Kadamba passes the food-safety checks before service. At Yuktāhār the first half hour of service projects the rest of the meal; from about 45 minutes before close staff check the counter every 30 minutes, open a second counter during surges, and at peak run two counters of handmade rotis cooked without oil (confirmed). At Kadamba the supervisor calls for more, ready in 10 to 15 minutes on site. At Bakul the manager phones Hafizpet, which cooks and sends it by road on an unknown lead time (single-sourced).
\item
  \textbf{Access.} A personal QR code is scanned at the counter to receive a plate matching the registered diet. The code can be reset at any time and staff can reassign it if a student eats at a different hall. When the scanner fails, staff write down ID and name by hand (single-sourced).
\item
  \textbf{Late entry.} At Kadamba the door reopens for 5 to 10 minutes after close, and students arriving after 9:30 also enter by the back door and are served; on one biryani day 25 came after close (single-sourced).
\item
  \textbf{Walk-ins and extra portions.} A student without a registration may pay the spot rate from a separate, smaller allocation. Kadamba cooks about 30 extra portions beyond registrations, described by the operator as for outside diners such as night security and by an earlier recording as 30 to 40 portions for paying walk-ins. Whether these are one allocation or two is unconfirmed (single-sourced).
\item
  \textbf{Leftover reuse and disposal.} Yuktāhār reuses leftovers and weighs student plate waste and kitchen waste (confirmed). At Kadamba held food is reheated for later service (candidate) and waste goes in a dustbin emptied by a garbage contractor, with no composting (single-sourced).
\item
  \textbf{Resale.} Outside the portal, a student transfers a registration through Mess Cell for a negotiated payment.
\item
  \textbf{Feedback.} Student input arrives through an in-app per-meal rating, a public email address and student-press coverage. Operator complaints pass from the institution\textquotesingle s site team or email to the operator\textquotesingle s manager and then the kitchen, at Kadamba through a WhatsApp group (single-sourced).
\end{itemize}

\subsection{Resources}\label{resources}

\begin{itemize}
\tightlist
\item
  The portal and the scan-level record it produces of who ate and when. The institution holds this record, which measures the registration-to-turnout gap for every meal it scans (confirmed for Kadamba breakfast; assumed for other halls and meals).
\item
  The operators\textquotesingle{} records. Yuktāhār\textquotesingle s books are listed above; utility readings feed decisions too, for example a Wednesday electricity peak from grinding idli batter informs how long machines run (single-sourced). Yuktāhār tried a commercial ERP to simplify purchasing and predict the rush and dropped it, judging its books enough for the cost and its staff unsuited to digital tools; office staff digitise the books later (single-sourced). Kadamba holds per-dish gram standards, a live count watched during service, a waste figure, retained 72-hour samples and a daily water-quality log (single-sourced). No written record of attendance, consumption or waste is known at Bakul. The registration-to-turnout gap is therefore already measured locally, by the people who cook.
\item
  Four halls of differing capacity and origin: Kadamba (purpose-built, largest), Bakul (a converted warehouse), Palash and Yuktāhār.
\item
  Production capacity split by cuisine: on-site kitchens at Kadamba and Yuktāhār, and one off-site kitchen with a transport link serving Bakul and Palash, which for Bakul means the Hafizpet kitchen, insulated boxes on the road and a hot-water holding setup at the hall (single-sourced).
\item
  Ingredients chosen on a quality-and-cost stance at Yuktāhār, which has used rock salt and cold-pressed oil for about six months (single-sourced).
\item
  Kitchen equipment and storage, including a cold room on first-in, first-out labels at Yuktāhār (confirmed). Kadamba\textquotesingle s toaster is broken and not yet replaced, and an LPG shortage was handled by cutting quantities and buying some items outside (single-sourced).
\item
  The personal QR credential issued to each student.
\end{itemize}

\begin{figure}
\centering
\includegraphics[width=1\linewidth,height=\textheight,keepaspectratio,alt={Yuktāhār\textquotesingle s kitchen on the day of the visit: the registers laid out, the cold room, and staff sorting vegetables at the weighing station.}]{Assets/plate-kitchen.jpg}
\caption{Yuktāhār\textquotesingle s kitchen on the day of the visit: the registers laid out, the cold room, and staff sorting vegetables at the weighing station.}
\end{figure}

\subsection{Constraints}\label{constraints}

\begin{itemize}
\tightlist
\item
  Registration is automatic, billing follows registration, and cancellations are capped at five per meal type per month, each at least four days ahead (confirmed).
\item
  The exit closes four days ahead, but Kadamba buys dry goods a week ahead and all other perishables one to three days ahead, so only non-perishable stock is committed before the exit closes and the kitchen\textquotesingle s own supply chain does not need the four-day lead (candidate).
\item
  At Kadamba each item passes a cooking-temperature check, a retained sample, a display plate, a serving-temperature check and a tasting before service (single-sourced), so the first batch must be cooked and approved before any sign of the day\textquotesingle s turnout exists (candidate). At Bakul the first batch is fixed earlier still, when it leaves Hafizpet to arrive at about 7 AM (single-sourced).
\item
  Cooked food is usable for only three to four hours (single-sourced). Raw material held back can be cooked later or on another day; cooked surplus cannot be recovered beyond a few reuses. The binding constraint on waste is therefore the cooked first batch, set as a share of registrations (fixed at Kadamba, adjusted weekly at Bakul), not raw procurement (candidate).
\item
  Cooked quantity can rise during service through top-up; nothing lowers it once the first batch is cooked. (Revises the first submission, which said nothing lets same-day activity change what is cooked.)
\item
  A Bakul top-up needs a phone call, cooking at Hafizpet and a 20 to 40 minute drive, against 10 to 15 minutes on site at Kadamba (single-sourced). Cooking time is unstated, so the true lead time is unknown. Late in the window a top-up cannot arrive before close, so food already held is the only cover for late arrivals; the manager\textquotesingle s reserve of about 50 portions at 9:20, several times the 10 to 15 late arrivals he describes, reads as that cover (candidate).
\item
  Bakul keeps food for about 70 percent of registrations even on low weekends (single-sourced), so any change to its quantities meets that floor first.
\item
  Food cooked off site must be held hot from dispatch to eating, at about 60 to 72 °C in insulated boxes. This keeps temperature and can cost texture (single-sourced for aloo paratha; untested for other items).
\item
  Kadamba cannot hold excess stock (single-sourced). Since dry goods are bought a week ahead, the limit reads as applying to perishable and cooked stock (candidate).
\item
  Cutting first-batch surplus at Kadamba also cuts food that 30 to 50 outside diners and about half the operator\textquotesingle s staff rely on (candidate).
\item
  Walk-in pricing is set about 70 percent above registered pricing.
\item
  Vegetarian and non-vegetarian food is segregated by registration type at service and planned separately at Kadamba; Palash and Yuktāhār serve vegetarian food only.
\item
  At Yuktāhār records are on paper, staff are not digitally inclined, and waste measurement took hold only once it was made blame-free (single-sourced).
\item
  Operators control quantities, but not registration rules, billing rules or the menu (consistent across three operators).
\item
  Exams, events and holidays reach a kitchen only informally, through whoever raises them in the planning meeting (confirmed).
\item
  Mess hours and the class schedule are set independently by separate authorities with no coordination mechanism.
\item
  Authority over mess operations is split across the Mess Committee, Mess Office and Warden, with CFS and CDS named by the Kadamba operator as its line; how these relate and who heads CDS is unconfirmed. The academic administration has no formal role.
\item
  Record-level turnout exists only for Kadamba breakfast in April 2026 and Yuktāhār in September 2026. Anything said about Bakul, Palash, or lunch and dinner outside Yuktāhār is an assumption until records are obtained.
\end{itemize}

\section{Initial Problem Framing}\label{initial-problem-framing}

Most breakfast registrations go unused. At Kadamba in April 2026 about 63 percent of registered breakfasts were not eaten; at Yuktāhār in September about 72 percent. At Yuktāhār lunch the same system is mostly used, with 56 to 89 percent of registrations eaten. Students are charged for every registered breakfast whether they eat it or not.

The kitchens know registration overstates breakfast, so none cooks for every registration. Each cooks to its own discounted guess and tops up during service. Kadamba\textquotesingle s first batch is still far above what arrives, while arrivals pile into the last fifteen minutes and cause shortages at the close. Surplus is partly reused, partly eaten by outsiders and the operators\textquotesingle{} staff, at Kadamba possibly held and reheated into later service, and the rest is thrown away. Kadamba\textquotesingle s operator puts waste at 5 to 10 percent a day, without saying of what, which is hard to square with a first batch near twice the turnout. Bakul shows the same pattern on its manager\textquotesingle s own account. He estimates at most about 60 percent come, cooks for 70 to 80 percent, and keeps food for about 70 percent even on low weekends, because his rule is no shortage even if some food is wasted (single-sourced; no Bakul records). He attributes low turnout partly to students not waking and partly to taste and dish pairing. The first agrees with what students interviewed elsewhere most often say. The second, like any effect of food softened in transport, is a candidate influence untested with Bakul residents.

Students do not choose to book breakfast; they are booked, on campus or not. To avoid paying for a breakfast they will not eat, they can cancel it at most five times a month per meal and at least four days ahead, or mark it skipped, which returns nothing. A student who decides on the morning itself has no option. Students built their own workaround, reselling registrations, because it pays back more than the system\textquotesingle s tools do.

\textbf{The problem: breakfast registration does not reflect who is going to eat, by a margin of roughly two thirds, while the same system works for lunch at Yuktāhār, the one hall with lunch records. Why not, and what in the system causes that gap? And why does each kitchen have to guess its own conversion from registrations to plates, when the data needed to calibrate it is already being recorded, and when one operator learns from its own books, another applies a fixed ratio, and a third adjusts its ratio by judgement?}

\end{document}
